\section{Related work}
\label{sec:relwork}

\paragrapht{Learning compact 3D scene representations.} 
Decomposing scenes into geometric primitives~(e.g., meshes, polygons, superquadric primitives) has been extensively studied~\cite{tulsiani2017learning, paschalidou2019superquadrics, paschalidou2020learning}, but these methods typically require 3D data~(e.g., point clouds) as input. SuperDec~\cite{fedele2025superdec} recently proposed feed-forward decomposition into superquadric primitives, but relies on pre-computed 3D point clouds and requires iterative application of an external 3D instance segmentation model to identify instances. The differentiable blocks world (DBW)~\cite{monnier2023differentiable} learns superquadric parameters directly from multi-view images through photometric optimization, but can only model up to 10 primitives and object-centric scenes. In contrast, we aim to learn an efficient solution that can estimate compact scene representations in a feed-forward manner, only given unposed multi-view images captured in real-world scenarios. 

\paragrapht{Feed-forward 3D Gaussian splatting.} 
Recently, 3D Gaussian Splatting (3DGS)~\cite{kerbl20233d} has gained significant attention for 3D reconstruction and novel view synthesis. However, it requires dense multi-view captures with known poses and time-consuming per-scene optimization. To address these limitations, recent approaches~\cite{charatan2024pixelsplat, chen2024mvsplat, du2023learning, johari2022geonerf, wang2021ibrnet, xu2024murf, yu2021pixelnerf, zhang2025transplat} have explored generalizable feed-forward networks~\cite{ye2024no, hong2024unifying, huang2025no, hong2024pf3plat, jiang2025anysplat, smart2024splatt3r} that synthesize novel views from sparse inputs by learning priors from large-scale datasets or leveraging foundation models such as pointmap regression models~\cite{wang2024dust3r, leroy2024grounding, wang2025vggt}. Despite these advances, existing feed-forward models~\cite{ye2024no, huang2025no} typically predict one or multiple Gaussians per-pixel, resulting in millions of Gaussians for multiple views or high-resolution images. This strategy produces excessive redundant Gaussians, degrading performance and causing artifacts as input views increase~\cite{wang2024freesplat, wang2025volsplat}. When incorporating semantic features through 2D-to-3D lifting, the excessive Gaussians create significant computational overhead~\cite{lee2025cf3, burgess2006spatial}. Previous works~\cite{zhou2024feature, fan2024large, burgess2006spatial, qin2024langsplat} address this by compressing semantic features into lower dimensions using specialized autoencoders, but this causes information loss and sub-optimal scene understanding~\cite{li2025langsplatv2}.


\paragrapht{Towards compact 3D Gaussian splatting.}
To address the redundancy problem in per-pixel Gaussian estimation, recent works~\cite{wang2024freesplat, huang2025longsplat, wang2025zpressor, zhang2024gaussian, ziwen2025long,jiang2025anysplat} have attempted to mitigate this issue by reducing the number of Gaussians post-hoc. FreeSplat~\cite{wang2024freesplat} and LongSplat~\cite{huang2025longsplat} iteratively add pixel-wise Gaussians only where existing projections are insufficient. ZPressor~\cite{wang2025zpressor} and Long-LRM~\cite{ziwen2025long} employ token merging to reduce redundant Gaussians with similar features before decoding them to per-pixel Gaussians. Anysplat~\cite{jiang2025anysplat} voxelizes the Gaussians in 3D space. However, these approaches do not fundamentally address the input-view bias inherent in per-pixel processing. EvolSplat~\cite{miao2025evolsplat} and VolSplat~\cite{wang2025volsplat} employ global voxel representations but remain domain-constrained or limited by fixed resolutions. We instead propose using learnable tokens to generate compact global Gaussians guided by input features, producing only essential Gaussians for scene representation.

